\chapter{Chest Injuries}

\section*{Definition and Overview}
Chest injuries are wounds or trauma to the thorax that may involve the ribs, sternum, lungs, heart, or the major blood vessels within the chest. They range from minor bruises and single rib fractures to life-threatening injuries such as a collapsed lung (pneumothorax), blood collecting in the chest cavity (haemothorax), damage to the heart, or tears of the great vessels. Because the chest houses the organs essential to breathing and circulation, chest trauma is a common cause of serious injury and death after accidents. Early recognition, rapid resuscitation, and timely treatment are critical, and the management of severe chest injury follows the principles of advanced trauma care.

\section*{Epidemiology}
Chest injuries account for a substantial proportion of trauma admissions worldwide and are present in a large number of fatal accidents. They are most common in young adult men involved in motor vehicle crashes, falls from height, and assaults. Road traffic collisions remain the leading single cause, while penetrating injuries from knives or guns are more common in certain urban settings. Older adults may sustain serious chest injuries -- such as multiple rib fractures -- from relatively minor falls because of osteoporosis and other age-related changes. The burden of chest trauma falls disproportionately on low- and middle-income countries, where road safety and trauma services are often less developed.

\section*{Aetiology and Pathophysiology}
\begin{itemize}
    \item \textbf{Blunt trauma:} motor vehicle crashes, falls, and crush injuries compress or shear the chest wall, fracturing ribs and bruising the lungs or heart
    \item \textbf{Penetrating trauma:} stab wounds and gunshot wounds directly puncture the chest wall, lungs, heart, or blood vessels
    \item \textbf{Pneumothorax:} air leaks into the space around the lung, causing it to collapse and impairing breathing; a tension pneumothorax builds pressure and can compress the heart and other lung
    \item \textbf{Haemothorax:} blood collects in the chest cavity, usually from injured blood vessels or the lung surface, reducing the space available for breathing
    \item \textbf{Flail chest:} multiple adjacent rib fractures allow a segment of the chest wall to move paradoxically, undermining effective breathing
    \item \textbf{Pulmonary contusion:} bruising of the lung tissue impairs gas exchange even when the chest wall is intact
    \item \textbf{Injury to the heart and great vessels:} severe blows or penetrating wounds can damage the heart muscle, valves, or aorta, causing rapid blood loss or heart failure
\end{itemize}

\section*{Clinical Features}
\begin{itemize}
    \item Chest pain that worsens with deep breathing, coughing, or movement
    \item Tenderness, bruising, swelling, or deformity over the ribs or sternum
    \item Shortness of breath, rapid breathing, or shallow, laboured breathing
    \item Abnormal or paradoxical movement of a section of the chest wall
    \item Cough, sometimes with blood-stained sputum
    \item Reduced or absent breath sounds on one side, detected by a doctor listening to the chest
    \item Signs of shock: pale, sweaty skin, rapid pulse, low blood pressure, and dizziness
    \item In severe cases, confusion, cyanosis, or loss of consciousness
\end{itemize}

\section*{Differential Diagnosis}
\begin{itemize}
    \item Heart attack and other cardiac conditions that mimic traumatic chest pain
    \item Aortic dissection presenting as sudden tearing chest or back pain
    \item Pulmonary embolism in a person with chest pain and breathlessness
    \item Musculoskeletal chest wall pain, such as muscle strain or costochondritis, without injury
    \item Pneumonia or pleurisy causing pleuritic pain and fever
    \item Spontaneous pneumothorax occurring without any trauma
\end{itemize}

\section*{When to Seek Medical Care}
Any significant blow to the chest, a penetrating wound, or persistent chest pain and breathing difficulty after injury requires urgent assessment. Even minor-looking injuries can be accompanied by internal damage.

\medskip
\textbf{Call 000 (or local emergency services) for:}
\begin{itemize}
    \item Difficulty breathing, severe shortness of breath, or gasping
    \item Chest pain after a fall, crash, or blow that is severe or worsening
    \item Coughing up blood after chest trauma
    \item Signs of shock: pale, sweaty, faint, or a rapid pulse
    \item Any penetrating wound to the chest or neck
    \item Loss of consciousness or suspected spinal injury
\end{itemize}

\section*{Diagnosis and Investigations}
\begin{itemize}
    \item \textbf{Primary survey:} rapid assessment of airway, breathing, and circulation in the emergency setting, following advanced trauma life-support principles
    \item \textbf{Physical examination:} inspection of the chest wall, listening for breath sounds, and checking for bruising, deformity, and crepitus over fractured ribs
    \item \textbf{Chest X-ray:} the first-line imaging test for rib fractures, pneumothorax, and haemothorax
    \item \textbf{CT scan:} detailed imaging for complex injuries, internal bleeding, or suspected damage to the aorta and other vessels
    \item \textbf{Bedside ultrasound:} rapid detection of air or blood around the lung during resuscitation
    \item \textbf{Blood tests:} full blood count, blood gases, and cross-match for transfusion in severe trauma
    \item \textbf{Electrocardiogram:} to check the heart after blunt or penetrating injury
\end{itemize}

\section*{Management}
\begin{itemize}
    \item \textbf{Pain relief:} simple analgesics, with careful use of stronger medication so that breathing is not suppressed
    \item \textbf{Simple rib fractures:} usually managed at home or in hospital with analgesia, ice, and breathing exercises
    \item \textbf{Pneumothorax:} drainage with a needle or chest tube to re-expand the lung; a tension pneumothorax requires immediate needle decompression
    \item \textbf{Haemothorax:} chest tube drainage and, if bleeding continues, surgery to control the source
    \item \textbf{Surgery:} for penetrating injuries, damage to the heart or large vessels, or persistent internal bleeding
    \item \textbf{Oxygen and breathing support:} oxygen therapy and, in severe cases, mechanical ventilation for significant pulmonary contusion or flail chest
    \item \textbf{Physiotherapy:} breathing exercises and early mobilisation to prevent pneumonia and lung collapse
    \item \textbf{Monitoring:} observation in hospital for pain, breathing, and signs of delayed complications
\end{itemize}

\section*{Complications}
\begin{itemize}
    \item Pneumonia and atelectasis (lung collapse) from shallow breathing and poor coughing
    \item Persistent air leak or recurrent pneumothorax
    \item Empyema -- infected fluid in the chest cavity
    \item Pulmonary contusion progressing to respiratory failure
    \item Damage to the heart, great vessels, or aorta causing major blood loss and shock
    \item Chronic chest wall pain or deformity after healing
    \item Post-traumatic stress and other psychological effects after serious injury
\end{itemize}

\section*{Prognosis and Outcomes}
Most minor chest injuries, including one or two simple rib fractures, heal well within a few weeks with pain relief, rest, and a gradual return to activity. Severe injuries involving major blood loss, a collapsed lung, or damage to the heart and great vessels are more serious and may require intensive care. Many people with major chest trauma recover fully, but some are left with chronic chest wall pain, reduced lung function, or psychological distress. Outcomes are generally better when emergency care is delivered rapidly and followed by active rehabilitation, including physiotherapy and graded return to normal activity.

\section*{Prevention and Self-Care}
\begin{itemize}
    \item Wear seatbelts in vehicles and helmets when riding motorcycles or bicycles
    \item Follow road safety rules and never drive under the influence of alcohol or drugs
    \item Use appropriate protective equipment at work and in contact sport
    \item Secure children in age-appropriate car seats
    \item After a minor injury, rest, apply ice, and use simple pain relief as advised
    \item Practise deep-breathing exercises and keep the chest moving as healing allows
    \item Stop smoking, since it increases the risk of pneumonia and complications after chest injury
\end{itemize}

\section*{Sources and Further Reading}
The following reputable sources provide further information on chest injuries:
\begin{itemize}
    \item NHS. \textit{Health A--Z: broken ribs and chest trauma information.} https://www.nhs.uk/conditions/
    \item Mayo Clinic. \textit{First aid and injury information.} https://www.mayoclinic.org/
    \item MSD Manual. \textit{Professional edition: thoracic trauma.} https://www.msdmanuals.com/
    \item MedlinePlus. \textit{Health topics on injuries and wounds.} https://medlineplus.gov/
    \item World Health Organization. \textit{Road traffic injury prevention.} https://www.who.int/
\end{itemize}
